% Author contribution, 21 November 2026.
% GPT-6 Astra Pro, /root/sixfolds; inherited reasoning effort not exposed.
\section{An exact nonsplit Gaussian sixfold source}
\label{sec:e8-exact-source}

\paragraph{Provenance and scope.}
The integer matrices below were taken as candidate data from
B.\ Severn, \emph{Weil classes on abelian sixfolds of Weil type for
\(\mathbb Q(i)\) with discriminant \(-3\)}, Hexagon 2610.00477v1,
Appendix A, printed p.40. No theorem or numerical assertion of that
draft is an input here. The group, the actual tuple used by the
reconstructed differential equation, its parabolic quotient, its cup
form, and the integral abelian factor are verified below. The period
calculation is HD-G17, with the existence calculation HD-I17.
The supplier calculation is a separate task: none of the statements
below asserts algebraicity of the Weil plane merely from the period rank.

\paragraph{HD-S22 (the actual tuple and its primitive abelian factor).}
Put \(K=\mathbb Q(i)\), \(\mathcal O=\mathbb Z[i]\), and let \(C_8\)
be the symmetric matrix with diagonal 2 and entries \(-1\) on the
edges
\[
 (1,3),(3,4),(4,5),(5,6),(6,7),(7,8),(2,4).
\]
Set \(H_0=C_8^{-1}\). Its entries are integers, \(\det C_8=1\),
and its leading principal minors are positive. The following matrices
act on columns, in the fundamental-weight basis:
\begingroup\small
\[
a=\begin{pmatrix}
0&0&1&0&0&0&0&0\\
1&1&1&1&1&1&0&0\\
1&1&1&2&2&2&2&1\\
-1&-1&-1&-1&-1&-1&-1&0\\
-1&-1&-2&-3&-3&-2&-1&-1\\
0&0&0&0&1&0&0&0\\
1&1&2&3&2&2&1&0\\
1&2&2&3&2&1&1&1
\end{pmatrix},\quad
b=\begin{pmatrix}
-1&-2&-3&-5&-4&-3&-2&-1\\
2&3&3&5&4&3&2&1\\
2&2&4&5&4&3&2&1\\
-2&-2&-3&-4&-4&-3&-2&-1\\
0&0&0&0&1&0&0&0\\
0&0&0&0&0&1&0&0\\
0&0&0&0&0&0&1&0\\
0&0&0&0&0&0&0&1
\end{pmatrix}.
\]
\endgroup
The matrix \(c\) equals the identity except that \(c_{11}=-1\)
and \(c_{31}=1\); put \(d=(abc)^{-1}\).
The actual positive-orientation tuple used by HD-I17 and HD-G17 is
\begin{equation}\label{eq:e8-actual-tuple}
 \boldsymbol g=
 \bigl(ia^5,
 a^{-4}ba^4,a^{-4}ca^4,
 a^{-3}ba^3,a^{-3}ca^3,
 a^{-2}ba^2,a^{-2}ca^2,
 a^{-1}ba,a^{-1}ca,b,c,-id^5\bigr).
\end{equation}
It has product one and generates
\[
 G=W(E_8)\cup iW(E_8),\qquad
 |G|=1393459200.
\]
Here the group identification is an exact finite permutation computation,
not an inference from the apparent form of the matrices. Starting with
the simple roots \(C_8u_j\), the simple reflections
\(1-C_8u_ju_j^t\) generate an orbit of exactly 240 roots of norm 2.
The action on these roots is faithful because they span the vector space.
Schreier--Sims gives the successive basic orbit lengths
\[
 240,126,32,15,8,3,2
\]
for the simple-reflection group. The matrices \(a,b,c\) belong to that
group and generate a group of the same order, \(696729600\).
For the extension, use the 480 distinct vectors \(R\cup iR\);
the corresponding first orbit length is 480 and the remaining lengths
are unchanged. The first eleven matrices of
\eqref{eq:e8-actual-tuple} generate the same subgroup as the conjugate
by \(a^{-4}\) of the first eleven matrices of the Appendix A tuple.
This is equality of explicitly displayed generating sets, up to order;
it does not claim a Hurwitz equivalence between tuples. Their order is
\(1393459200\), proving the asserted generation. The twelfth matrix
is already the inverse product of the preceding ones.
All these assertions, including membership in the simple-reflection
group and product one, are checked with exact integer or Gaussian
rational arithmetic by
\texttt{e8\_group\_and\_tuple.py} and
\texttt{e8\_actual\_parabolic.py}.

The matrices preserve \(h_0(x,y)=\bar x^tH_0y\).
The eight simple-root lines and the connected root diagram show that
the commutant is scalar: a commuting matrix preserves each root line,
and adjacent root lines force their scalars to agree. Equivalently,
products of the rank-one operators \((1-s_j)/2\) along paths in the
diagram span all root outer products and therefore all of \(M_8(K)\).
The representation is absolutely irreducible. Its character field is
exactly \(K\): all entries are Gaussian integers, and the central
element \(z=iI_8\) has character value \(8i\).

The branch codimensions are \((4,1^{10},8)\). The eigenvalue angles,
in units of \(1/12\), of the two non-reflections are
\[
 (0,0,0,0,6,6,6,6),\qquad (1,2,4,5,8,9,9,10).
\]
Each of the ten reflections has one angle \(1/2\) and seven zero
angles. Thus both the sum of the angles and the sum of the nonzero
complementary angles are 11. The parabolic dimension is
\(4+10+8-2\cdot8=6\), and the two Hodge dimensions are
\(11-8=3\). One published input with a complete proof is
Dettweiler--Wewers, \emph{Variation of parabolic cohomology and
Poincar\'e duality}, S\'eminaires \& Congr\`es 13 (2006),
Theorem 2.3, pp.155--157. Their proof applies Riemann--Roch to the
two Deligne extensions and uses the exact parabolic dimension to turn
both lower bounds into equalities. Here the local system is unitary
and has no invariants, as required.

For twelve distinct ordered branch points, gluing copies of a cut sphere
with these sheet permutations gives a connected topological cover.
The charts \(t\mapsto t^e\) at a point of inertia order \(e\)
give its compact Riemann surface \(C\); thus it is a smooth projective
complex curve. Riemann--Hurwitz gives
\[
 g(C)=1+\frac{53|G|}{24}=3077222401.
\]
Use the regular fibre \(G\), with monodromy acting on the right and
left deck maps \(L_h:g\mapsto hg\). These deck maps form a genuine
left action. The action of \(\mathbb Z[G]\) on \(J(C)=\operatorname{Pic}^0(C)\)
used throughout is the covariant action
\[
 a_h=(L_h)_*=(L_{h^{-1}})^*,
\]
where the first star denotes pushforward of divisors (equivalently
the norm on line bundles), and the last star denotes pullback of line
bundles. On curve cohomology put \(T_h=(L_{h^{-1}})^*\). Then \(T\)
is a genuine representation, Poincar\'e duality satisfies
\(a_{h,*}\operatorname{PD}=\operatorname{PD}T_h\), and the Abel map
identifies \(a_h^*\) on \(H^1(J(C))\) with \(T_{h^{-1}}\) on
\(H^1(C)\). The possible base-point translation in the Abel map
acts trivially on cohomology. These identities fix the group-ring
action, rather than identifying it with an opposite action implicitly.

Take the root \(v=C_8u_1\), of norm 2, and its orthogonal projection
\(p(x)=v\,h_0(v,x)/2\). Define
\begin{equation}\label{eq:e8-integral-projector}
 e=\frac8{|G|}\sum_{g\in G}
       \operatorname{Re}h_0(v,\rho(g^{-1})v)\,g,
 \qquad D=\frac{|G|}{8}=174182400.
\end{equation}
Then \(e^2=e=e^*\), \(e\mathbb Q[G]e\simeq K\), and
\(De\in\mathbb Z[G]\). Here \(*\) sends \(g\) to \(g^{-1}\).
For completeness, write
\(f(g)=h_0(v,\rho(g^{-1})v)/2=\operatorname{tr}(p\rho(g^{-1}))\).
Finite-group averaging gives matrix-coefficient orthogonality:
the average of \(\rho(g)A\rho(g)^{-1}\) is
\(\operatorname{tr}(A)I_8/8\), since its commutant is scalar.
Applied to matrix units, this proves that
\((8/|G|)\sum f(g)g\) is a rank-one matrix idempotent in the
\(K\)-block. Its conjugate is orthogonal to it: multiplying a summation
index by \(z\) multiplies a crossed coefficient by \(-1\).
Their sum is exactly \eqref{eq:e8-integral-projector}; the corner is
the conjugate pair of scalar corners, hence \(K\) over \(\mathbb Q\).
Also \(f(g^{-1})=\overline{f(g)}\), proving self-adjointness.
Finally \(H_0,v,\rho(g)\) are integral over \(\mathcal O\), so the
real parts in \eqref{eq:e8-integral-projector} are integers.

Consequently
\[
 B=\operatorname{im}\bigl(De:J(C)\longrightarrow J(C)\bigr)^0
\]
is an actual abelian subvariety, not merely an isogeny symbol.
It has dimension six. The central automorphism \(i_B=a_z^{-1}|_B\) makes
\(\mathcal O\) act integrally on \(B\); the restricted theta
polarization has Rosati involution complex conjugation. The preceding
Hodge calculation makes this a Weil sixfold of signature \((3,3)\).

\paragraph{HD-S23 (geometric cup form and integral source bounds).}
Here is the exact convention comparison used in the determinant
calculation. For the column tuple \(g_1,\ldots,g_{12}\), set
\(P_i=g_1\cdots g_{i-1}\). A parabolic cocycle is represented by
\[
 c_i=(1-g_i)a_i,\qquad \sum_iP_ic_i=0,
\]
modulo the tuples \(((g_i-1)x)_i\). Put
\(F_i=\sum_{j<i}P_jc_j\). For two cocycles define
\begin{equation}\label{eq:e8-cup-formula}
 S(u,v)=\sum_i h_0(P_i^{-1}F_i^u+a_i^u,c_i^v),\qquad H=iS.
\end{equation}
This is the geometric parabolic cup form. To check the precise map,
use the right-module convention of Dettweiler--Wewers,
Theorem 2.6, pp.158--160. Reverse the indices and set
\[
 G_j=g_{13-j}^t,\quad v_j=c_{13-j}^t,\quad
 v'_j=-a_{13-j}^t,\quad v_j^*=\bar c_{13-j}^{,t}H_0.
\]
The right-module relation is the transpose of \(\sum P_ic_i=0\),
and its coboundaries are the transposes of the displayed coboundaries.
The dual right matrix is \(G_j^*=g_{13-j}^{-1}\), since
\(\bar g_i^tH_0=H_0g_i^{-1}\). These are row/column and path-order
conventions, with an explicit isomorphism of the cochain quotients;
no conjugacy between \(g_i^t\) and \(g_i^{-1}\) is used.
Their formula becomes
\[
 -\sum_i h_0(c_i^u,a_i^v)
 -\sum_{i<l}h_0(P_i^{-1}P_lc_l^u,c_i^v).
\]
The cocycle relation replaces the second sum by
\(\sum_i h_0(P_i^{-1}F_i^u+c_i^u,c_i^v)\).
Unitarity and \(c_i=(1-g_i)a_i\) give
\[
 h_0(c_i^u,c_i^v)=h_0(a_i^u,c_i^v)+h_0(c_i^u,a_i^v),
\]
which gives exactly \eqref{eq:e8-cup-formula}.
Changing a preimage \(a_i\) by a fixed vector changes no term,
because \(\ker(1-g_i)\perp\operatorname{im}(1-g_i)\).
Moreover the identity
\[
 S(u,v)+\overline{S(v,u)}
 =\sum_i\bigl(h_0(F_{i+1}^u,F_{i+1}^v)
                  -h_0(F_i^u,F_i^v)\bigr)=0
\]
shows skew-Hermitian symmetry. The published cup comparison also
proves descent through coboundaries; it is checked independently on
the full cocycle matrix below.

For reproducibility, take the first independent columns of \(1-g_i\)
as a matrix \(V_i\), and write \(c_i=V_iu_i\). The relation matrix
\([P_iV_i]_i\) has rank 8 on a 22-dimensional space. Its kernel has
dimension 14. The coboundary matrix has rank 8 and lies in that kernel.
Column reduction, adjoining the coboundaries first, gives an integral
\(22\times6\) quotient basis \(Q\). Its entries are provided in
\texttt{e8\_actual\_parabolic.json}. Computing
\eqref{eq:e8-cup-formula} on the full 14-dimensional cocycle space
gives rank 6 and has the entire coboundary space as its radical.
On \(Q\), the resulting Hermitian matrix is
\begin{equation}\label{eq:e8-actual-H}
 H=\begin{pmatrix}
0&i&-1-i&-1&-2&1+i\\
-i&0&-1-2i&-1-i&-2+2i&1\\
-1+i&-1+2i&0&2i&0&i\\
-1&-1+i&-2i&0&-2+2i&1+i\\
-2&-2-2i&0&-2-2i&0&-2-2i\\
1-i&1&-i&1-i&-2+2i&2
\end{pmatrix},\qquad \det H=-24.
\end{equation}
Thus the ordinary determinant class is \([-3]\), since
\(8=N_{K/\mathbb Q}(2+2i)\). This determinant was computed on the
actual tuple \eqref{eq:e8-actual-tuple}; the old Appendix A tuple has
a separately computed Gram determinant \(-48\), which is not used
to identify the present component. Twenty independently constructed
contiguous-return logarithms preserve \eqref{eq:e8-actual-H}.
That is a consistency check, not the method used to identify the form.

We give the actual comparison to \(B\), including its rational scalar.
On a fibre of the covering, regarded as functions on \(G\), define
\[
 j_x(g)=\frac{h_0(v,\rho(g)x)}2.
\]
Right monodromy sends \(j_x\) to \(j_{\rho(k)x}\). To distinguish
the other index, write \(j_{v,x}(g)=h_0(v,\rho(g)x)/2\). The
cohomological deck representation above satisfies
\[
 T_hj_{v,x}=j_{\rho(h)v,x}.
\]
This depends conjugate-linearly on \(v\), so it is the conjugate
representation \(\bar\rho\) in that index. In particular
\(T_zj_x=-ij_x\).
Matrix-coefficient orthogonality gives
\[
 \sum_{g\in G}\overline{j_x(g)}j_y(g)
       =\frac{|G|}{16}h_0(x,y).
\]
The image is the \(T_z=-i\) eigensummand of \(e\) on the regular
local system. This equality supplies a left inverse and hence an
isomorphism on its parabolic cohomology summand. After extending over
the branch points and applying Poincar\'e duality, define
\[
 \Phi(u)=\operatorname{PD}\bigl(j_u+\overline{j_u}\bigr)
       \in eH_1(C,\mathbb Q)=H_1(B,\mathbb Q).
\]
Here conjugation is coefficientwise and is applied to the conjugate
local-system class. Explicitly,
\[
 \Phi(iu)=\operatorname{PD}(ij_u-i\bar j_u)
        =\operatorname{PD}T_{z^{-1}}(j_u+\bar j_u)
        =a_{z^{-1},*}\Phi(u).
\]
Thus \(\Phi\) is \(K\)-linear for \(i_B=a_z^{-1}\). On
\(H^1(B)\), however, pullback by \(i_B\) is \(T_z\): its
\(\rho\) block has eigenvalue \(+i\), and the coefficient block
just constructed is the conjugate \(-i\) block. This explicitly
identifies which of the two Weil lines the supplier identity uses.
The two equal-character cup terms vanish by the central deck action.
For the convention
\(h_B(x,y)=E_B(x,i_By)+iE_B(x,y)\), expansion of the other two
terms gives
\begin{equation}\label{eq:e8-source-gram}
 h_B(\Phi u,\Phi v)
   =2i\int_C\overline{j_u}\smile j_v
   =\frac{|G|}{8}\,H(u,v)=D H(u,v).
\end{equation}
This supplies the polarization comparison by an explicit cohomological
map. In particular, multiplying the rank-six Gram determinant by
\(D^6=N(D^3)\) does not alter its norm class.

There is also a useful integral consequence. The basis \(Q\) above is
integral, and all matrices \(\rho(g)\) are Gaussian integral.
Therefore the pulled-back cocycle \(j_u\), for \(u\in\mathcal O^6\),
has coefficients in \(\tfrac12\mathcal O\); its sum with its conjugate
is integral. The parabolic condition kills the loops that are capped
at the branch points, so this is an integral class on the compact curve.
The same holds after multiplication by \(i\). Consequently, in the
\(\Phi\)-coordinates and with \(H_s=DH\),
\begin{equation}\label{eq:e8-source-lattice-sandwich}
 \mathcal O^6\ \subseteq\ \Lambda_B\ \subseteq\ H_s^{-1}\mathcal O^6.
\end{equation}
The second containment is integral duality: the restricted alternating
pairing and its translate by \(i_B\) are integral, so
\(h_B(\Lambda_B,\mathcal O^6)\subseteq\mathcal O\).

An independent bound uses the self-adjoint idempotent directly. Write
\(L=H_1(C,\mathbb Z)\), a unimodular symplectic lattice, and
\(\Lambda_B=L\cap eL_{\mathbb Q}\). This intersection is primitive.
Restriction \(L^*\to\Lambda_B^*\) is surjective, because
\(L/\Lambda_B\) is torsion-free. Under the unimodular pairing and
the self-adjoint projection this identifies \(eL\) with
\(\Lambda_B^*\). Thus
\[
 D\Lambda_B^*=DeL\subseteq L\cap eL_{\mathbb Q}=\Lambda_B.
\]
Every elementary divisor of the restricted polarization divides \(D\).
If \(\Delta_B\) is their product, then
\begin{equation}\label{eq:e8-source-product-bound}
 \Delta_B\mid D^6.
\end{equation}

Finally the actual homomorphism \(q=De:J(C)\to B\), composed with
the Abel map, supplies a curve with controlled theta degree.
The pullback form is \(q^*\theta_B=D^2e^*\theta_C\), and its trace
on the holomorphic tangent space is \(6D^2\).
For a direct verification of the trace-to-degree identity, choose
holomorphic differentials \(\omega_j\) on \(C\) orthonormal for
\((i/2)\int_C\omega_j\wedge\bar\omega_k\). The invariant theta
form of the Jacobian pulls back under the Abel map to
\((i/2)\sum_j\omega_j\wedge\bar\omega_j\). The self-adjoint
idempotent is an orthogonal projection of rank six on these forms,
so integrating its pullback gives its trace. Consequently
\begin{equation}\label{eq:e8-abel-degree}
 \int_C q^*\theta_B=6D^2=182037050818560000.
\end{equation}
This is the degree of the pushed-forward effective one-cycle, including
its generic mapping multiplicity. It does not replace the separate
construction of a detecting three-cycle.

\paragraph{HD-S24 (computable transport to each specified target lattice).}
Every rational Hermitian space over \(K=\mathbb Q(i)\) of signature
\((3,3)\) and determinant class \([-3]\) is explicitly isometric to
\[
 H_{\rm ref}=U\perp U\perp\langle1,-3\rangle,
 \qquad U=\begin{pmatrix}0&1\\1&0\end{pmatrix}.
\]
Here is a constructive proof using only the already cited form of
Meyer's theorem (Serre, \emph{A Course in Arithmetic}, Chapter IV,
Section 3.2, Corollary 2 to Theorem 8, p.43). Apply it to the underlying
indefinite rational quadratic forms of ranks 12 and then 8. At each
stage, an isotropic vector \(v\), a vector \(w\) with \(h(v,w)=1\),
and the replacement \(w-h(w,w)v/2\) split off an actual \(U\).
Diagonalize the remaining plane as \(\langle a,b\rangle\).
Its determinant condition gives \(ab=-3N(c)\) for some \(c\in K^\times\).
Scaling its second vector by \(a/c\) changes the plane to
\(a\langle1,-3\rangle\).

For every \(a\in\mathbb Q^\times\), the five-variable form
\[
 X^2+Y^2-3Z^2-3W^2-aT^2
\]
is indefinite, so Meyer supplies a nonzero rational zero.
Its coordinate \(T\) is nonzero: a zero with \(T=0\) would make
3 a norm from \(\mathbb Q(i)\), unless all coordinates vanished.
The norm impossibility follows by clearing denominators and reducing
a putative \(r^2+s^2=3t^2\) modulo 3, then infinite descent.
Put \(x=(X+iY)/T\) and \(y=(Z+iW)/T\). Then
\(N(x)-3N(y)=a\), and the explicit matrix
\begin{equation}\label{eq:e8-binary-transport}
 T_a=\begin{pmatrix}x&3\bar y\\y&\bar x\end{pmatrix}
 \quad\hbox{satisfies}\quad
 T_a^*\begin{pmatrix}1&0\\0&-3\end{pmatrix}T_a
     =a\begin{pmatrix}1&0\\0&-3\end{pmatrix},\qquad
 \det T_a=a.
\end{equation}
All searches in this construction are effective: enumerate rational
vectors by height until the proved isotropic or norm equation is met.
The orthogonal complements, diagonalization, and changes of basis are
exact linear algebra. This proves termination; it does not assert a
closed height bound in terms of the polarization product alone.

In particular, for a specified integral Gaussian target Gram matrix
\(H_t\), this procedure constructs \(F\in\operatorname{GL}_6(K)\)
with \(F^*H_tF=H_s\), where \(H_s=DH\) is the concrete source form
in \eqref{eq:e8-source-gram}. Let \(M\) be any positive integer
clearing all real and imaginary denominators of \(FH_s^{-1}\).
The sandwich \eqref{eq:e8-source-lattice-sandwich} gives
\[
 MF\Lambda_B\subseteq\mathcal O^6.
\]
Transporting a compatible complex structure by \(F\) therefore gives
the actual isogeny
\[
 \phi=MF:B\longrightarrow A_t,\qquad
 \phi^*\theta_t=M^2\theta_B.
\]
For an effective three-cycle \(Z\), the projection formula gives
\(\deg_{\theta_t}\phi_*Z=M^6\deg_{\theta_B}Z\).
Moreover \(\phi_*\) and \(\phi^*\) preserve the Weil summands,
and \(\phi_*\phi^*=\deg(\phi)\) on cohomology. Thus this actual
pushforward preserves nonvanishing of a Weil projection.
For a target specified by an arbitrary order-stable rank-12 integral
lattice, express \(FH_s^{-1}\mathcal O^6\) in its given rational
\(\mathbb Z\)-basis and clear those denominators instead. The same
isogeny and degree formula apply; no unproved conductor estimate is
inserted.

There is one completely numerical reference case. Exact congruence
reduces \eqref{eq:e8-actual-H} to
\(U^2\perp\begin{psmallmatrix}6&6+6i\\6-6i&8\end{psmallmatrix}\).
The columns of the actual congruence matrix are
\[
 T=\begin{pmatrix}
 1&0&-i&1+3i&-4-4i&-8i\\
 0&-i&-1+i&-2i&5i&-4+6i\\
 0&0&0&-i&i&2i\\
 0&0&0&0&1&0\\
 0&0&0&0&0&1\\
 0&0&1&-1&3-2i&6
 \end{pmatrix},\qquad \det T=1.
\]
The second block diagonalizes to \(\langle6,-4\rangle\); scaling
its second vector by \((3+3i)/2\) gives \(6\langle1,-3\rangle\).
The matrix \(\begin{psmallmatrix}3&3\\1&3\end{psmallmatrix}\)
then gives the reference form by \eqref{eq:e8-binary-transport}.
In formulas, put
\[
 C=\begin{pmatrix}1&-1-i\\0&1\end{pmatrix}
       \operatorname{diag}(1,(3+3i)/2),\qquad
 R=\begin{pmatrix}3&3\\1&3\end{pmatrix},\qquad
 P=T\operatorname{diag}(I_4,CR^{-1}).
\]
Then \(P^*HP=H_{\rm ref}\).
For the additional scalar \(D\), use
\[
 x=17280,\qquad y=2880+5760i,
 \qquad N(x)-3N(y)=D,
\]
and \(\operatorname{diag}(D,1)\) on each \(U\)-block.
Set \(S_D=\operatorname{diag}(D,1,D,1,T_D)\), with \(T_D\)
from \eqref{eq:e8-binary-transport}, and \(F=S_DP^{-1}\).
The resulting exact isometry \(F\) has
\(D F H_s^{-1}\in M_6(\mathcal O)\).
Thus \(M=D\) works for the fixed target lattice \(H_{\rm ref}\),
whose polarization product is 3. The matrices and this denominator
are checked in \texttt{e8\_primitive\_and\_transport.py/.json}.

\paragraph{Proof audit.}
Author checks, 21 November 2026, GPT-6 Astra Pro, inherited effort:
HD-S22 was read in full after the exact finite-group recomputation;
HD-S23 was read in full after the direct actual-tuple cup computation,
including the row/column convention map, the coefficient pairing scalar,
the integral lattice sandwich, and the idempotent dual-lattice argument;
HD-S24 was read in full, including termination of the rational searches
and both lattice inclusions defining the isogeny. These are author
proof checks, not independent reproofs. The exact arithmetic programs
separately certify the displayed group orders, ranks, Gram determinant,
source constants, and reference isometry. External whole-proof reviews
are recorded as separate events when completed. The unresolved supplier
task remains separate from these proved geometric and arithmetic facts.
